% ============================================================================
% Appendix insert: place immediately after the proof of Lemma C.8, before the
% existing ``Proof of Theorem 2(i)'' (or after it if preferred).  Needs amsmath,
% amssymb, amsthm; a lemma environment and proof environment must exist.
% Also needs bibliography items for Carbery--Wright and Gaussian
% hypercontractivity (see theorem2_references.bib).
% ============================================================================

\subsubsection*{Quantitative rule credit under Gaussian initialization}
\label{app:t2-gaussian}

We refine \cref{lem:leading} from generic nonvanishing to a quantitative
high-probability bound.  Write
\begin{equation}
 F_n(\varepsilon)=
 \left.\frac{\partial\Phi_{\mathrm{rule}}}{\partial V_{ab}}
 \right|_{\theta=\varepsilon\widetilde\theta}
 =\varepsilon^6P_n(\widetilde\theta)+r_n(\varepsilon),
 \qquad
 |r_n(\varepsilon)|\le\frac{C_{\mathrm{rem}}}{n}
 \left(\frac{\varepsilon^7}{\varepsilon_0}+\varepsilon^{10}\right).
 \label{eq:quant-rule-expansion}
\end{equation}
We may weaken this to $C_{\mathrm{rem}}(\varepsilon^7/(n\varepsilon_0)+\varepsilon^{10})$
when using \cref{eq:t2-gaussian-epsilon}. The remainder estimate holds whenever the deterministic hypotheses of
\cref{lem:leading} are satisfied and
$\varepsilon\le\varepsilon_0$,
$C_{\mathrm{att}}\varepsilon^4\le1/2$.

\begin{lemma}[A square-free monomial controls Gaussian $L^2$]
\label{lem:quant-rule-L2}
Let $P$ be a homogeneous polynomial of degree six in independent centered
Gaussian coordinates, and suppose the coefficient of the square-free
monomial $X_1\cdots X_6$ is $c\ne0$, with
$X_i\sim\mathcal N(0,\sigma_i^2)$ and $\sigma_i>0$.
Then
\begin{equation}
  \|P\|_{L^2}\ge |c|\prod_{i=1}^6\sigma_i.
  \label{eq:quant-rule-L2}
\end{equation}
\end{lemma}
\begin{proof}
Put $Y=\prod_{i=1}^6(X_i/\sigma_i)$.  Independence and centering give
$\mathbb E Y^2=1$.  In $\mathbb E[PY]$, a nonzero contribution from a
monomial of $P$ requires each of $X_1,\ldots,X_6$ to occur to a positive odd
power: otherwise at least one centered Gaussian has an odd total exponent.
Since every monomial of $P$ has total degree six, the only surviving one is
$X_1\cdots X_6$.  Thus
$\mathbb E[PY]=c\prod_{i=1}^6\sigma_i$.
The claim follows by Cauchy--Schwarz.
\end{proof}

\begin{lemma}[Degree-six Gaussian anti-concentration]
\label{lem:quant-rule-AC}
There is an absolute constant $C_6>0$ such that, for any nonzero real
polynomial $P$ of total degree at most six evaluated on an independent
Gaussian vector with nondegenerate coordinate variances and every
$0<\rho<1$,
\begin{equation}
 \mathbb P\!\left\{
   |P|\ge\left(\frac{\rho}{C_6}\right)^6\|P\|_{L^2}
 \right\}\ge 1-\rho.
 \label{eq:quant-rule-AC}
\end{equation}
\end{lemma}
\begin{proof}
After rescaling coordinates, it suffices to work with a standard Gaussian
vector.  The Carbery--Wright inequality for Gaussian measure gives, for a
polynomial of degree at most six and every $t>0$,
\begin{equation}
 \mathbb P\{|P|\le t\}
 \le 6C_{\mathrm{CW}}
       \left(\frac{t}{\|P\|_{L^1}}\right)^{1/6}
 \label{eq:quant-rule-CW}
\end{equation}
with a universal $C_{\mathrm{CW}}$; see
\cite{carbery2001distributional}.
Gaussian hypercontractivity
\cite{janson1997gaussian} implies
$\|P\|_{L^4}\le 3^3\|P\|_{L^2}=27\|P\|_{L^2}$.
Interpolation gives
\begin{equation}
 \|P\|_{L^2}\le
 \|P\|_{L^1}^{1/3}\|P\|_{L^4}^{2/3},
 \qquad
 \|P\|_{L^1}\ge 3^{-6}\|P\|_{L^2}.
\end{equation}
Substitution in \eqref{eq:quant-rule-CW} yields
\begin{equation}
 \mathbb P\{|P|\le t\}
 \le C_6\left(\frac{t}{\|P\|_{L^2}}\right)^{1/6},
 \qquad C_6:=\max\{1,18C_{\mathrm{CW}}\}.
\end{equation}
Take $t=(\rho/C_6)^6\|P\|_{L^2}$.
\end{proof}

\begin{lemma}[High-probability lower bound for the actual rule gradient]
\label{lem:quant-rule-gradient}
Fix $n\ge2$, $K\ge2$, $0<p<1/2$, and $0<\delta<1$.
Assume the architecture and uniform-start task of \cref{lem:leading}, and suppose its $P_n$ is a degree-six polynomial in
independent centered Gaussian attention/embedding entries.  Assume the six
distinguished entries have standard deviations $\sigma_1,\ldots,\sigma_6$.
Let $\mathcal E_{\mathrm{geom}}$ be an event, of probability at least
$1-\delta/2$, on which the deterministic hypotheses of \cref{lem:leading} hold with
fixed nonrandom $\varepsilon_0,C_{\mathrm{att}},C_{\mathrm{rem}}>0$.
Define $\kappa_{p,K}$ and $A_{n,\delta}$ by
\eqref{eq:t2-gaussian-threshold}.  If $\varepsilon$ satisfies
\eqref{eq:t2-gaussian-epsilon}, then
\begin{equation}
 \mathbb P\!\left\{
  \|\nabla_\theta\Phi_{\mathrm{rule}}
     (\varepsilon\widetilde\theta)\|
       \ge\tfrac12\varepsilon^6 A_{n,\delta}
 \right\}\ge1-\delta.
 \label{eq:quant-rule-gradient}
\end{equation}
\end{lemma}
\begin{proof}
\cref{lem:leading} gives a square-free degree-six monomial whose coefficient
$c_n$ satisfies
\begin{equation}
 |c_n|\ge
 \frac{p(1-p)\min\{1-2p,1/4\}}{3Kn\sqrt{d_h}}
 =\frac{\kappa_{p,K}}{n\sqrt{d_h}}.
 \label{eq:quant-rule-coeff}
\end{equation}
The six distinct entries, for the selected block-2 head and readout row
$a=s^\ast+1$, are
\[
 \tilde E_{s^\ast,0},\quad \tilde E_{w^\star,1},\quad
 \tilde E_{w^\star,2},\quad
 (\tilde{\mathbf W}_Q)_{0,0},\quad
 (\tilde{\mathbf W}_K)_{0,1},\quad
 (\tilde{\mathbf W}_V)_{b,2}.
\]
They belong to different matrix entries even when their numerical indices
coincide. The polynomial contains no readout or MLP-bias coordinate.
By Lemma~\ref{lem:quant-rule-L2},
$\|P_n\|_{L^2}\ge
 \kappa_{p,K}(n\sqrt{d_h})^{-1}\prod_i\sigma_i$.
Applying Lemma~\ref{lem:quant-rule-AC} with $\rho=\delta/2$ gives
\begin{equation}
 \mathbb P\{|P_n|\ge A_{n,\delta}\}\ge1-\delta/2.
\end{equation}
The union bound (which requires no independence between the signal and
geometry events) gives
\begin{equation}
 \mathbb P\!\left(
  \{|P_n|\ge A_{n,\delta}\}\cap\mathcal E_{\mathrm{geom}}
 \right)\ge1-\delta.
\end{equation}
On this intersection, the deterministic remainder estimate
\eqref{eq:quant-rule-expansion} and the third and fourth bounds in
\eqref{eq:t2-gaussian-epsilon} imply
\begin{align}
 \frac{|r_n(\varepsilon)|}{\varepsilon^6}
 &\le C_{\mathrm{rem}}\left(
       \frac{\varepsilon}{n\varepsilon_0}+\varepsilon^4\right)
 \le\frac{A_{n,\delta}}2.
\end{align}
Therefore
\begin{equation}
 \|\nabla_\theta\Phi_{\mathrm{rule}}\|
 \ge |F_n(\varepsilon)|
 \ge\varepsilon^6 |P_n|-|r_n(\varepsilon)|
 \ge\tfrac12\varepsilon^6A_{n,\delta}.
\end{equation}
\end{proof}

\paragraph{Binary specialization and scope.}
For $K=2$ and $p=1/4$, one has $\kappa_{p,K}=1/128$;
hence the threshold in \eqref{eq:quant-rule-gradient} is
\begin{equation}
 \frac{\varepsilon^6}{256n\sqrt{d_h}}
 \biggl(\prod_{i=1}^6\sigma_i\biggr)
 \biggl(\frac{\delta}{2C_6}\biggr)^6.
\end{equation}
If $n$, $d_h$, all $\sigma_i^{-1}$, $\varepsilon_0^{-1}$ and the
uniform remainder constants are polynomially bounded in a size parameter,
then, for inverse-polynomial $\delta$ and provided the geometry event has
the stated probability, both an admissible $\varepsilon$ and the resulting
positive lower bound can be inverse-polynomial in that parameter.
The statement does not prove the requisite geometric event for arbitrary
widths or arbitrary Gaussian scaling, nor does it show that unrestricted
optimization dynamics amplify this gradient.
\paragraph{An explicit Gaussian-weight geometry event.}
\label{app:t2-gaussian-geometry}
The conditional event can be verified for the following specified law.
Fix all widths with $d_{\mathrm{res}}\ge\max\{K,3\}$. Draw each normalized
matrix/embedding block $M_b\in\mathbb R^{r_b\times c_b}$ independently
with independent entries $N(0,\tau_b^2)$, where every $\tau_b>0$ is
specified before sampling. A zero readout may instead be fixed. Let $B$
be the number of Gaussian blocks, and let value-token embedding rows
have common entry standard deviation $\tau_E$. Set
\[
 a_0=\max_b 2\tau_b\sqrt{2[(r_b+c_b)\log9+\log(8B/\delta)]},\qquad
 c_{\mathrm{emb}}=\frac{\delta\tau_E}{4K^{3/2}}.
\]
Use normalized MLP biases with $1\le|\tilde b_m|\le2$ and fix
$N_b=2\sqrt{w_{\max}}$, where $w_{\max}$ is the largest MLP width.
All relative biases are fixed initially at zero. Then, with probability
at least $1-\delta/2$, every block has operator norm at most $a_0$ and
$\sigma_{\min}(E_{\mathrm{val}})\ge c_{\mathrm{emb}}$. Thus the geometric
conditions of \cref{ass:init} have fixed, nonrandom envelopes; choose
$\varepsilon_0$ below its threshold and those of \cref{lem:leading}.
The resulting $C_{\mathrm{att}},C_{\mathrm{rem}}$ are the deterministic
constants for these envelopes, not values fitted after observing a seed.

For the operator-norm assertion, take $1/4$-nets on both unit spheres,
each of cardinality at most $9^{\mathrm{dimension}}$. The norm is at most
twice the maximal bilinear form on these nets. Gaussian tails and a union
bound give failure probability at most $\delta/(4B)$ for each displayed
block bound. For embedding conditioning, let $D_i$ be the distance of
row $i$ to the span of the other $K-1$ rows. Conditional on those rows,
at least one orthogonal projection is $N(0,\tau_E^2)$, so
$\Pr(D_i<t)\le\sqrt{2/\pi}\,t/\tau_E$. For any unit combination of rows,
choose a coefficient of magnitude at least $K^{-1/2}$ and project
orthogonally to the span of the other rows; hence
$\sigma_{\min}(E_{\mathrm{val}})\ge\min_i D_i/\sqrt K$.
Taking $t=\delta\tau_E/(4K)$ bounds conditioning failure by $\delta/4$.
The two failure bounds sum to at most $\delta/2$.

Fan-in Gaussian weights are covered by taking $\tau_b=c_b^{-1/2}$ and
embedding variance $\tau_E^2=d_{\mathrm{res}}^{-1}$. This explicit witness
still uses non-Gaussian, bounded-away-from-zero biases. It does not cover
ordinary zero-bias or fully Gaussian-bias initialization. Nor does it
assert that all envelope-dependent expansion constants are width-uniform.
For the formal-regime bias law $\tilde b_m=\operatorname{sign}(Z_m)(1+|Z_m|)$
in \nolinkurl{handcoded/regime_both.py}, with independent standard-normal
$Z_m$ and $N_{\mathrm{bias}}$ such coordinates, instead use
\[
 N_b=\sqrt{w_{\max}}\left(1+\sqrt{2\log(16N_{\mathrm{bias}}/\delta)}\right),
\]
and replace $8B/\delta$ above by $16B/\delta$. Matrix norms, bias norms,
and embedding conditioning then fail with total probability at most
$\delta/8+\delta/8+\delta/4=\delta/2$.
